\begin{abstract}
Models trained on educational data may amplify inequalities that conventional metrics miss. This paper presents AuditEdu-XAI, a Design Science Research artifact, with a free public interface, that audits, on the test set only, XGBoost classifiers of three ENEM editions with Disparate Impact, group error rates and TreeSHAP. All non-white groups and incomes up to 2.5 minimum wages violate the four-fifths rule in all three editions, and the model amplifies the race/color disparity in the labels in 14 of 15 group-edition combinations (13 with a 95\% CI entirely below zero). Flipping 50\% of the positive training labels of \textit{pardo} candidates, with a negative control, brought the group's Disparate Impact to near zero; the ROC AUC changed, on average, by at most 0.032.
\end{abstract}

\begin{resumo}
Modelos treinados com dados educacionais podem amplificar desigualdades invisíveis às métricas usuais. Apresenta-se o AuditEdu-XAI, artefato de \textit{Design Science Research}, com interface pública e gratuita, que audita, só no conjunto de teste, classificadores XGBoost de três edições do ENEM com Impacto Disparate, taxas de erro por grupo e TreeSHAP. Todos os grupos não brancos e rendas de até 2,5 salários mínimos violam a regra dos quatro quintos nas três edições, e o modelo amplifica a disparidade de raça/cor dos rótulos em 14 de 15 combinações grupo-edição (13 com IC95\% inteiramente negativo). Rebaixar 50\% dos rótulos positivos de treino dos pardos, com controle negativo, quase zerou seu Impacto Disparate; a AUC-ROC variou, em média, até 0,032.
\end{resumo}
